\documentclass[11pt,a4paper]{article}

\usepackage[utf8]{inputenc}
\usepackage[T1]{fontenc}
\usepackage{lmodern}
\usepackage[margin=22mm,headheight=15pt]{geometry}
\usepackage{microtype}
\usepackage{parskip}
\usepackage{enumitem}
\usepackage{booktabs}
\usepackage{longtable}
\usepackage{tabularx}
\usepackage{xcolor}
\usepackage{hyperref}
\usepackage{fancyhdr}
\usepackage{listings}
\usepackage{amsmath}
\usepackage{amssymb}

\definecolor{AionBlue}{HTML}{126E82}
\definecolor{AionNavy}{HTML}{17324D}
\definecolor{AionLight}{HTML}{EAF5F7}
\definecolor{CodeBg}{HTML}{F4F6F7}
\definecolor{WarningBg}{HTML}{FFF4DD}
\definecolor{RiskRed}{HTML}{A63A3A}

\hypersetup{
  colorlinks=true,
  linkcolor=AionNavy,
  urlcolor=AionBlue,
  pdftitle={AION Native Communication and Discourse Repair},
  pdfauthor={Tessaris / AION Engineering}
}

\pagestyle{fancy}
\fancyhf{}
\lhead{AION Native Communication and Discourse Repair}
\rhead{Technical Record --- 14 October 2026}
\cfoot{\thepage}

\setlist[itemize]{leftmargin=*,itemsep=2pt,topsep=3pt}
\setlist[enumerate]{leftmargin=*,itemsep=2pt,topsep=3pt}
\renewcommand{\arraystretch}{1.18}

\lstdefinestyle{aioncode}{
  backgroundcolor=\color{CodeBg},
  basicstyle=\ttfamily\small,
  breaklines=true,
  frame=single,
  rulecolor=\color{AionLight},
  columns=fullflexible,
  keepspaces=true,
  showstringspaces=false
}
\lstset{style=aioncode}

\newcommand{\statusbox}[1]{%
  \begin{center}
  \fcolorbox{AionBlue}{AionLight}{\parbox{0.92\textwidth}{#1}}
  \end{center}}

\newcommand{\warningbox}[1]{%
  \begin{center}
  \fcolorbox{RiskRed}{WarningBg}{\parbox{0.92\textwidth}{#1}}
  \end{center}}

\title{\textcolor{AionNavy}{\Huge AION Native Communication and Discourse Repair}\\[4mm]
\Large Contract Safety, Multi-turn Dialogue, Live Evidence Migration and the Route to Open Conversation}
\author{Tessaris / AION Engineering}
\date{Technical snapshot: 14 October 2026}

\begin{document}
\maketitle

\begin{abstract}
This record documents the work performed after the initial connected-English system was
introduced.  It covers the Native Discourse and Dialogue Engine, the Native Communication
Academy, strict source and answer-contract validation, removal of contaminated historical
evidence, correction-queue control, delayed retention, cross-subject scheduling, runtime
services, SQLite storage on the TESSARISQ38 SD volume, verification tests and the next native
conversation architecture.  The central system boundary remains unchanged: Qwen may prepare
sealed teaching material, but it is absent while AION interprets, reasons and answers.
\end{abstract}

\statusbox{\textbf{Capability statement.} AION now demonstrates genuine but bounded English
comprehension.  It can form multi-sentence situation models, retrieve evidence, answer several
semantic question families, construct connected responses and participate in structured
multi-turn examinations without a language model in the student path.  It is not yet an
unrestricted conversational system.  Open-ended intent recognition, robust reference
resolution, topic changes, pragmatic language and broadly varied native expression remain
active engineering goals.}

\tableofcontents
\newpage

\section{Scope and Relationship to Earlier Documentation}

This document supplements, rather than replaces, the earlier technical records:

\begin{itemize}
  \item \path{docs/latex/AION_Native_English_Intelligence_Architecture.tex}; and
  \item \path{docs/latex/AION_Connected_English_Advancement_Section_2026-10-11.tex}.
\end{itemize}

Those records describe lexical memory, the semantic spider, Word Wall exercises, the Practice
Gym, reading, semantic transfer and the first discourse implementation.  The present record
focuses on the later repairs needed to make connected-English evidence trustworthy and to move
AION from short structured answers toward sustained native communication.

\section{Non-negotiable Student--Teacher Boundary}

The architecture deliberately separates teaching from performance.

\begin{lstlisting}
Qwen/local teacher
    -> proposes a sealed situation, question and verified answer contract
    -> disconnects

AION/native student
    -> parses the passage and question
    -> builds or updates a situation model
    -> retrieves lexical and acquired concept evidence
    -> selects and reasons over propositions
    -> constructs an English response
    -> checks support, ambiguity and confidence

Sealed grader
    -> reveals the expected meaning only after AION has answered
    -> records success, correction or delayed retest
\end{lstlisting}

The following invariants apply:

\begin{enumerate}
  \item Qwen must not supply, rewrite or polish AION's student response.
  \item The expected answer is withheld until after the native response exists.
  \item A correct answer counts only when supported by the represented source evidence.
  \item Invalid teacher material is rejected before AION sees it.
  \item Practice rewards guide scheduling but do not, by themselves, prove mastery.
  \item AION may request clarification or state insufficient evidence instead of guessing.
\end{enumerate}

\section{Native Discourse and Dialogue Engine}

The principal implementation is:

\begin{center}
\path{backend/modules/aion_cognition/native_discourse_dialogue.py}
\end{center}

\subsection{Situation model}

The engine converts connected text into an explicit representation containing:

\begin{itemize}
  \item entities and their textual mentions;
  \item events, agents, actions and affected objects;
  \item locations and explicit time expressions;
  \item clause-level negation;
  \item causal links and event order;
  \item pronoun and reference relationships;
  \item evidence sentences and source identifiers; and
  \item unresolved or ambiguous relationships.
\end{itemize}

The response path is proposition-based rather than answer-string retrieval:

\[
\text{question}
\rightarrow \text{intent}
\rightarrow \text{selected event/relation}
\rightarrow \text{evidence}
\rightarrow \text{propositions}
\rightarrow \text{surface response}.
\]

Current sealed examination families include entity, action, complete object, location, time,
cause, coreference, first event, next event and contradiction.  Definition exercises are not
reported as discourse comprehension.

\subsection{Evidence and confidence}

Each response plan contains the native answer, requested intent, propositions, evidence,
confidence and support checks.  A response is not considered fully successful merely because
its text resembles the answer.  It must also be grounded in the current situation model and the
sealed answer contract.

\section{Native Communication Academy}

The multi-turn curriculum is implemented in:

\begin{center}
\path{backend/modules/aion_cognition/native_communication_academy.py}
\end{center}

The academy advances through the following evidence ladder:

\begin{longtable}{p{0.28\textwidth}p{0.12\textwidth}p{0.50\textwidth}}
\toprule
Stage & Turns & Required capability \\
\midrule
One-sentence answer & 1 & One supported native answer \\
Connected explanation & 2 & Answer plus connected evidence \\
Short paragraph & 3 & Several propositions in a coherent response \\
Passage discussion & 4 & Repeated reference to one represented passage \\
Five-turn dialogue & 5 & Persistent state across a short conversation \\
Ten-turn dialogue & 10 & Complete, varied and evidence-backed conversation \\
Cross-subject dialogue & 10 & Conversation grounded in different source subjects \\
Delayed dialogue retest & 10 & Retained performance after scheduled delay \\
\bottomrule
\end{longtable}

Graduation requires two complete passes from distinct passages.  Advanced stages additionally
require source diversity.  An earlier implementation silently accepted fewer turns when a
passage did not provide enough contracts; those receipts are now retained for audit but excluded
from capability claims.

\section{Contract-Safety Repair}

The most important recent discovery was that high-volume activity did not necessarily represent
valid English testing.  Some imported pages contained indexes, headings or catalogue fragments.
The parser then created questions with outputs such as pronouns, particles, truncated phrases or
long lists.  Examples included prompts equivalent to ``the event involving I'' and object keys
such as ``me'' or incomplete clause fragments.

\subsection{Connected-prose gate}

The \texttt{\_passage\_is\_prose} gate now rejects material that exhibits index-like structure,
including:

\begin{itemize}
  \item too few words or no connected sentence/clause structure;
  \item excessive title-cased fragments;
  \item semicolon-delimited lists dominated by very short entries;
  \item contents, index, appendix, bibliography or glossary headings;
  \item repeated figure, table, volume, section or page-number references; and
  \item excessively large candidate windows.
\end{itemize}

This filter is applied while reading candidates are assembled and again before a sealed contract
can enter an examination.

\subsection{Complete object phrases}

The \texttt{\_object\_answer\_safe} gate now requires a bounded, complete noun phrase.  It rejects:

\begin{itemize}
  \item standalone pronouns;
  \item particles and prepositions;
  \item connective or function-word endings;
  \item parser spill-over from adjacent clauses; and
  \item phrases without a lexically supported noun head.
\end{itemize}

Valid examples include ``the damaged sensor'' and ``the loose cable.''  Invalid examples include
``me,'' ``up the curtain'' and ``the fact that only.''

Object questions are now created only when the event has a grounded subject, a usable action and
a unique subject--action match within the passage.  The compiled question is anchored to that
event, for example:

\begin{lstlisting}
What complete object did Leona inspect?
\end{lstlisting}

This replaces weak questions such as ``What was affected in the event involving I?''

\subsection{Other role gates}

Additional validation requires:

\begin{itemize}
  \item grounded named entities for agent questions;
  \item lexically place-like heads for location phrases;
  \item non-auxiliary, non-connective action surfaces;
  \item bounded answers for all families; and
  \item grounded event descriptions for sequence answers.
\end{itemize}

\section{Reference Ambiguity and the Current Ten-turn Finding}

The first repaired ten-turn examination achieved nine correct turns out of ten.  The missed item
asked what ``they'' referred to in a combined passage that contained two separate occurrences of
``they.''  AION answered ``Gold,'' while the sealed contract expected ``Ore.''  Because the
question did not identify which occurrence was intended, the item was ambiguous.

This is not evidence that reference resolution is complete.  It identifies the next contract
repair: every coreference question must carry a mention identifier, sentence index or sufficient
local quotation to select one occurrence unambiguously.  Future reference contracts should use
a form such as:

\begin{lstlisting}
In the sentence beginning "Ore minerals concentrated...",
what does the first occurrence of "they" refer to?
\end{lstlisting}

The mention identity must remain available to the native planner even when Qwen rephrases the
surface question.

\section{Dialogue Diversity and Anti-scaffolding Controls}

Several controls now prevent volume from masquerading as transfer:

\begin{enumerate}
  \item No more than two questions from one semantic family may fill a dialogue session.
  \item A ten-turn session therefore requires several different abilities rather than ten object
        lookups.
  \item An identical passage signature is not reused within the same curriculum stage.
  \item Cross-subject sessions combine three passages from three distinct subject families.
  \item Exact passage signatures, source identities and family counts are retained as evidence.
  \item Qwen proposals that expose the answer, change the semantic role or select an unknown
        contract are rejected.
  \item Native compiled contracts fill safe empty teacher slots so an invalid Qwen draft cannot
        starve learning.
\end{enumerate}

The distinction is important:

\[
\text{repeated structured success} \neq \text{unfamiliar language transfer}.
\]

Only source-disjoint, sealed and model-free examinations contribute to the stronger transfer
claims.

\section{Corrections, Retests and Queue Control}

Every failed valid contract creates a verified remedial obligation.  The correction stores the
expected meaning and an explanation, but the later retest runs closed-book.  Delays are scheduled
at approximately one day, seven days and thirty days.

Recent queue repairs include:

\begin{itemize}
  \item verified corrections are processed before new sessions are generated;
  \item new work pauses while a large due queue is being drained;
  \item a failed session creates at most four immediate remedial obligations;
  \item duplicate failures from the same semantic family do not create repeated obligations;
  \item repeated failures are delayed rather than hammered continuously; and
  \item retests sharing a passage reuse one in-memory situation model.
\end{itemize}

The model cache is particularly important for ten-turn retention.  Without it, the same passage
could be parsed ten times during one retest batch.  The optimized path parses it once and executes
all associated native plans against the same immutable model.

\section{Historical Evidence Migration}

The stricter contract rules were applied to the existing SD-card ledger.  No records were
deleted.  Invalid sessions and retests were assigned explicit invalidation statuses and excluded
from ability metrics.

\begin{longtable}{p{0.46\textwidth}r p{0.34\textwidth}}
\toprule
Audit classification & Count & Treatment \\
\midrule
Invalid contract quality & 1,026 sessions & Preserved; excluded from evidence \\
Incomplete stage & 21 sessions & Preserved; excluded from graduation \\
Invalid teacher entity contract & 1 session & Preserved; excluded from evidence \\
Invalid contract retests & 5,902 & Cancelled as student obligations \\
Incomplete-stage retests & 55 & Cancelled as student obligations \\
\bottomrule
\end{longtable}

The apparent fall in session totals was therefore a correction of the evidence ledger, not a loss
of acquired lexical memory.  Immediately after migration, valid communication accuracy rose and
new repaired sessions began passing.  The migration marker is stored in
\texttt{native\_communication\_meta}, preventing the expensive full audit from rerunning at every
dashboard request.

\section{Persistence and SQLite Tables}

The live learning database remains:

\begin{center}
\path{/Volumes/TESSARISQ38/TessarisLearning/AION/native_lexical_intelligence.sqlite3}
\end{center}

Relevant discourse tables include:

\begin{itemize}
  \item \texttt{native\_discourse\_models};
  \item \texttt{native\_discourse\_entities};
  \item \texttt{native\_discourse\_events};
  \item \texttt{native\_dialogue\_sessions};
  \item \texttt{native\_dialogue\_turns};
  \item \texttt{native\_dialogue\_corrections};
  \item \texttt{native\_discourse\_exams}; and
  \item \texttt{native\_discourse\_exam\_remediation}.
\end{itemize}

Communication Academy tables include:

\begin{itemize}
  \item \texttt{native\_communication\_sessions};
  \item \texttt{native\_communication\_turns};
  \item \texttt{native\_communication\_retests}; and
  \item \texttt{native\_communication\_meta}.
\end{itemize}

Session rows retain stage, source, subject family, passage signature, passage, target and completed
turn counts, correct and evidence-supported turns, teacher identity, student-model-use flag,
timestamp and validity status.  Turn rows retain the family, question, sealed expected answer,
native answer, correctness, confidence, propositions, evidence and claim checks.

SQLite connections are governed through:

\begin{center}
\path{backend/modules/aion_cognition/sqlite_connection.py}
\end{center}

The shared connection helper applies the configured busy timeout and the project's single-writer
discipline, reducing the SD-card lock collisions seen in earlier learning builds.

\section{Runtime Services and Dashboard}

The relevant live processes are:

\begin{longtable}{p{0.34\textwidth}p{0.56\textwidth}}
\toprule
Service & Responsibility \\
\midrule
\texttt{aion\_native\_reading\_service.py} & Continues book ingestion, reading, contextual evidence and retention \\
\texttt{aion\_mastery\_curriculum\_service.py} & Schedules curriculum, transfer, discourse and communication work \\
\texttt{uvicorn backend.desktop\_app:app} & Serves the local Learning Centre and evidence APIs \\
\bottomrule
\end{longtable}

The Learning Centre endpoint is:

\begin{center}
\url{http://127.0.0.1:8080/api/aion/english-intelligence/panel}
\end{center}

Progress snapshots are cached briefly so browser polling does not repeatedly scan large SD-card
tables.  The API publishes completed sections progressively and exposes component refresh errors;
a temporarily cached panel must not be confused with stalled learning.

\section{Dated Live Evidence Snapshot}

The following measurements were observed on 14 October 2026 after the repair and evidence
migration.  They are operational evidence, not timeless specifications.

\begin{longtable}{p{0.55\textwidth}r}
\toprule
Measure & Value \\
\midrule
Valid Communication Academy sessions & 20 \\
Valid sessions passed & 13 \\
Valid dialogue turns & 72 \\
Overall valid communication accuracy & 86.1\% \\
Recent communication accuracy & 95.8\% \\
Current communication stage & Ten-turn dialogue \\
First repaired ten-turn result & 9/10 \\
Discourse examinations & 4,716 \\
Recent discourse accuracy & 97.3\% \\
Reading sentences processed & 109,642 \\
Reading sentences pending & 1,155 \\
Recent reading accuracy & 91.5\% \\
Recent contextual-vocabulary accuracy & 97.7\% \\
Concepts verified in reading context & 4,890 \\
Recalled-use concepts & 3,250 \\
Contextual-use evidence, lifetime & 166 \\
Currently strict contextual-stage concepts & 4 \\
Productive-use concepts & 162 \\
Recent Practice Gym attempts & 502 at 92.0\% \\
Recent Language Workshop attempts & 30 at 100.0\% \\
\bottomrule
\end{longtable}

\warningbox{The recent free-transfer result was 100\%, but it contained only five attempts.
It is encouraging, not statistically sufficient.  Likewise, structured ten-turn performance
does not yet demonstrate arbitrary open conversation.}

\section{What AION Can and Cannot Currently Do}

\subsection{Demonstrated capabilities}

AION can currently:

\begin{itemize}
  \item interpret many ordinary sentence structures;
  \item represent agents, actions, objects, locations, time, negation and causes;
  \item preserve a bounded situation across connected sentences;
  \item answer sealed questions from represented evidence;
  \item construct short connected explanations and paragraphs;
  \item run structured multi-turn dialogue without Qwen answering for it;
  \item report uncertainty when its native plan lacks the requested role; and
  \item retain corrections through scheduled closed-book retests.
\end{itemize}

\subsection{Capability boundary}

The present evidence does not yet prove reliable handling of:

\begin{itemize}
  \item arbitrary open-ended user questions;
  \item informal speech, humour, irony and implied social meaning;
  \item multiple identical pronouns in a long combined passage;
  \item long conversations with topic changes and interruptions;
  \item varied natural expression across arbitrary subjects; or
  \item questions whose relevant knowledge is absent from AION's represented memories.
\end{itemize}

The correct description is therefore \emph{genuine early, bounded English comprehension and
structured conversation}, not fluent general-purpose conversation.

\section{Next Build: Native Conversation Runtime}

Current vocabulary and reading should continue, but additional dictionary volume is no longer the
main engineering constraint.  The next system should convert connected meaning into robust open
interaction.  This is not principally a waiting problem.  More spider pages, book sentences and
definition exercises will strengthen lexical coverage, but they cannot by themselves create the
state management, intent interpretation, reference resolution and response planning required for
free conversation.  Those mechanisms must be implemented and examined directly.

\subsection{System objective}

The Native Conversation Runtime should accept an arbitrary owner utterance and produce a native,
evidence-governed response without using next-token model generation in the student path.  Its
top-level flow is:

\begin{lstlisting}
owner utterance
  -> token, clause and discourse parse
  -> mention and reference resolution
  -> open intent/request frame
  -> dialogue-state query
  -> lexical, episodic and knowledge retrieval
  -> reasoning and proposition selection
  -> response plan
  -> native grammatical realisation
  -> evidence, ambiguity and confidence checks
  -> transactional dialogue-state update
\end{lstlisting}

The runtime has two valid outcomes: a supported response, or an explicit clarification/unknown
response.  An unsupported guess is not a valid third outcome.

\subsection{Proposed implementation boundaries}

The following modules are proposed; these paths describe the intended separation of
responsibilities and do not imply that all files already exist:

\begin{longtable}{p{0.48\textwidth}p{0.42\textwidth}}
\toprule
Proposed module & Responsibility \\
\midrule
\path{native_reference_graph.py} & Mention identities, candidate antecedents, reference scoring and ambiguity \\
\path{native_open_intent.py} & Speech acts, requested roles, scope, assumptions and missing information \\
\path{native_dialogue_memory.py} & Persistent turn state, topics, beliefs, corrections and unresolved obligations \\
\path{native_knowledge_retrieval.py} & Evidence retrieval from acquired concepts, reading and episodic memories \\
\path{native_response_realizer.py} & Proposition ordering, clause planning, morphology and surface generation \\
\path{native_conversation_runtime.py} & End-to-end orchestration, policy checks and transactional commit \\
\path{native_open_conversation_academy.py} & Sealed training conversations, grading, corrections and retention \\
\bottomrule
\end{longtable}

Existing discourse and communication modules should remain the tested foundation.  The new
runtime should call their stable parsing and evidence interfaces rather than copy their logic.

\subsection{Mention-specific reference graph}

Every pronoun or referring phrase requires a stable mention identifier with sentence position,
candidate antecedents, compatibility features and confidence.  The engine must ask for
clarification when two antecedents remain plausible.

A mention record should contain at least:

\begin{itemize}
  \item a stable mention identifier and canonical entity or event identifier;
  \item turn, sentence and token span;
  \item surface text and mention type;
  \item grammatical number, person and available semantic class;
  \item discourse role, recency and current-topic salience;
  \item candidate antecedents with positive and negative evidence; and
  \item selected antecedent, confidence and resolution status.
\end{itemize}

Candidate resolution should be an explicit scored decision rather than a nearest-name shortcut.
One possible bounded score is:

\[
S(a,m)=w_rR+w_gG+w_sS+w_tT+w_cC-w_kK,
\]

where $R$ is recency, $G$ grammatical compatibility, $S$ semantic compatibility, $T$ topic
salience, $C$ clause-role compatibility and $K$ contradiction or conflict evidence.  Resolution
is permitted only when the best candidate exceeds both an absolute threshold and the second-best
candidate by a configured margin.  Otherwise the runtime generates a clarification request.

Identical surface pronouns must remain separate mention nodes.  For example, two occurrences of
``they'' in different source sentences cannot share one answer contract merely because their text
is identical.  Sealed examinations must identify the intended mention without exposing its
antecedent.

\subsection{Open intent interpretation}

The fixed question-family classifier should be extended into a compositional request frame:

\begin{lstlisting}
user utterance
  -> speech act / intent
  -> requested entity, event, property or relation
  -> time and discourse scope
  -> presuppositions
  -> missing or ambiguous information
\end{lstlisting}

This frame must support questions, instructions, corrections, comparisons, explanations,
clarification and changes of subject.

The request frame should be serialisable and contain:

\begin{lstlisting}
{
  "speech_act": "question | instruction | correction | assertion | social",
  "requested_relation": "entity | event | property | cause | comparison | ...",
  "arguments": [{"role": "target", "entity_id": "..."}],
  "time_scope": {"start": null, "end": null, "relative_to": "..."},
  "discourse_scope": ["current_turn", "previous_turn", "active_topic"],
  "presuppositions": [],
  "ambiguities": [],
  "missing_requirements": [],
  "confidence": 0.0
}
\end{lstlisting}

Interpretation should combine lexical concepts, clause structure, dialogue state and current
topic.  It must not depend on a finite list of complete question strings.  Existing semantic
families can remain typed primitives inside the larger compositional frame.

If a question contains two plausible targets, the runtime should preserve both hypotheses and ask
which one the owner intends.  It must not write a guessed interpretation into persistent memory.

\subsection{Persistent dialogue state}

The native state should persist:

\begin{itemize}
  \item people, objects and events introduced by the user;
  \item accepted facts and their provenance;
  \item current and previous topics;
  \item unresolved references and questions;
  \item earlier AION claims and confidence; and
  \item user corrections and belief revisions.
\end{itemize}

State updates should be transactional and evidence-labelled.  A correction must supersede the
affected proposition without silently deleting its history.

Dialogue state should be divided into four layers:

\begin{enumerate}
  \item \textbf{Turn state}: the current utterance, parse hypotheses, intent and unresolved spans.
  \item \textbf{Situation state}: active people, objects, events, relations, time and location.
  \item \textbf{Discourse state}: current topic, topic stack, pending questions, references and
        conversational obligations.
  \item \textbf{Belief state}: accepted propositions, provenance, confidence, contradictions and
        superseding corrections.
\end{enumerate}

Every turn should follow a prepare--validate--commit pattern:

\begin{lstlisting}
prepare candidate state changes
  -> verify evidence and references
  -> detect contradiction or ambiguity
  -> produce response
  -> commit only supported changes
  -> retain superseded facts as audit history
\end{lstlisting}

Topic changes should push the previous topic onto a recoverable stack.  Phrases such as ``going
back to the inspection'' should reactivate the relevant topic and its entities without treating
them as newly introduced facts.

Proposed persistence tables include
\texttt{native\_conversation\_sessions},
\texttt{native\_conversation\_turns},
\texttt{native\_conversation\_mentions},
\texttt{native\_conversation\_topics},
\texttt{native\_conversation\_propositions},
\texttt{native\_conversation\_corrections} and
\texttt{native\_conversation\_exams}.  All should reside in the same governed SD-card database
unless profiling proves that a separate derived index is required.

\subsection{Searchable, structured knowledge}

English competence and world knowledge must remain distinguishable.  AION cannot answer an
arbitrary question merely because it can parse its grammar.  Reading output therefore needs a
searchable proposition index rather than retrieval of whole books or raw dictionary entries.

Each durable knowledge item should contain:

\begin{itemize}
  \item subject, predicate and object/property identifiers;
  \item time, location, negation and modality where available;
  \item source, sentence span and acquisition method;
  \item confidence and verification status;
  \item links to supporting or contradicting propositions; and
  \item the acquired lexical concepts used to interpret it.
\end{itemize}

Retrieval should begin with the request frame, expand through acquired semantic relations and
return a small evidence bundle.  The response planner should not receive an entire book when three
propositions answer the question.  Source text may be fetched for verification, but it must not
replace the structured representation.

\subsection{Native language realiser}

The answer generator should operate over intended propositions:

\[
\text{reasoning result}
\rightarrow \text{ordered propositions}
\rightarrow \text{clause plan}
\rightarrow \text{tense/reference plan}
\rightarrow \text{word forms}
\rightarrow \text{surface response}
\rightarrow \text{support check}.
\]

It should produce multiple valid forms without memorising one answer template.  Grading should
compare semantic frames, grammatical constraints and supported claims rather than exact strings.

The realiser should operate in bounded stages:

\begin{enumerate}
  \item select only propositions relevant to the response goal;
  \item order propositions by answer, evidence, qualification and conclusion;
  \item choose sentence and clause boundaries;
  \item assign subject, action, object, complements and connectors;
  \item select tense, aspect, number, determiners and pronouns;
  \item realise stored or productive morphology;
  \item check agreement, reference integrity and unsupported additions; and
  \item select one of several meaning-equivalent surface forms.
\end{enumerate}

Variation should come from explicit grammatical alternatives, not random word substitution.  For
example, active/passive choice, clause combination, connective selection and evidence placement
may vary only when the proposition graph proves that meaning is preserved.

Before emission, each clause should carry a support map from surface claim back to proposition and
evidence identifiers.  A clause lacking support is removed, qualified as inference, or replaced by
an uncertainty statement.

\subsection{Pragmatic behaviour}

The runtime must learn to request clarification, distinguish fact from inference, report missing
knowledge, recognise likely non-literal language, explain evidence and revise an earlier answer.
These behaviours are part of English competence, not optional presentation polish.

At minimum, the policy layer needs explicit outcomes for:

\begin{longtable}{p{0.28\textwidth}p{0.62\textwidth}}
\toprule
Condition & Native behaviour \\
\midrule
Ambiguous reference & Name the alternatives and ask which was intended \\
Missing evidence & State what is unknown and what information would resolve it \\
Fact versus inference & Label the inference and cite its supporting propositions \\
Contradictory evidence & Present the conflict and avoid silently choosing one source \\
Owner correction & Record provenance, supersede the affected belief and acknowledge the update \\
Non-literal cue & Preserve a literal hypothesis and a likely pragmatic hypothesis until resolved \\
Low confidence & Give a calibrated answer or clarification rather than a confident guess \\
Topic change & Save the old topic, establish the new one and retain a route back \\
\bottomrule
\end{longtable}

Humour and implication should not initially be treated as unconstrained creativity.  They should
be represented as competing literal and intended-meaning hypotheses, learned from sealed examples
and graded on whether AION requests clarification when evidence is weak.

\subsection{Turn-level safety and observability}

Every native turn should expose an internal receipt containing intent, selected evidence,
propositions, mention resolutions, confidence, ambiguities, unsupported-claim count, state changes
and model-use flags.  The user-facing response need not display the receipt by default, but the
Learning Centre and examination ledger must retain it for audit.

The runtime must fail closed if a teacher-model process remains attached to the student response
path.  A turn with an uncertain model-use boundary is ineligible for native capability evidence.

\section{Open Conversation Academy}

Qwen may create sealed conversations at scale, subject to the same student boundary.  Training
should progress through:

\begin{lstlisting}
one supported answer
  -> two-sentence explanation
  -> native paragraph
  -> passage discussion
  -> five-turn dialogue
  -> ten-turn dialogue
  -> twenty-turn dialogue with topic changes
  -> cross-subject dialogue
  -> delayed source-disjoint retest
\end{lstlisting}

Teacher scenarios should deliberately include unfamiliar names, paraphrases, informal wording,
ambiguous requests, corrections, disagreements, missing information and changes of topic.  The
teacher output is validated and sealed before AION starts.

\subsection{Teacher contract}

For each session, Qwen may produce:

\begin{itemize}
  \item the private world state and source propositions;
  \item a sequence of owner utterances and allowed topic changes;
  \item intended mentions and semantic request frames;
  \item meaning-equivalent answer criteria;
  \item forbidden unsupported claims;
  \item expected clarification points; and
  \item verified corrections and later retest contracts.
\end{itemize}

The material must pass a deterministic validator before the answer key is sealed.  Qwen then
disconnects.  AION receives only the visible conversation.  The grader opens the private contract
after each native response and records semantic correctness, evidence support, grammar, reference
integrity, state retention and uncertainty calibration separately.

\subsection{Adaptive progression}

Training should allocate work according to the failure type:

\begin{itemize}
  \item a reference error returns to mention discrimination and later uses a new surface form;
  \item an intent error receives a simpler request-frame contrast before open re-examination;
  \item a knowledge-retrieval error receives evidence-selection practice, not a grammar drill;
  \item a realisation error keeps the correct proposition plan and retrains only surface grammar;
  \item an unsupported claim receives evidence-bound response practice; and
  \item a state-loss error receives a shorter conversation with the same memory operation.
\end{itemize}

One verified correction should be followed by a simpler closed-book retry, then an unfamiliar
free-transfer test, then delayed source-disjoint retention.  Repeating the failed question with the
answer visible is not learning evidence.

\subsection{Curriculum mix}

An initial open-conversation allocation should be approximately:

\begin{itemize}
  \item 25\% reference resolution and dialogue-state continuity;
  \item 20\% arbitrary intent and clarification;
  \item 20\% knowledge retrieval and evidence selection;
  \item 20\% native response planning and realisation;
  \item 10\% correction, contradiction and belief revision; and
  \item 5\% implication, humour and other pragmatic language.
\end{itemize}

The mix should adapt as evidence accumulates.  Perfect elementary performance should reduce its
own scheduling weight; it must not crowd out the ability that remains weak.

\section{Acceptance Criteria for Conversational English}

The system should not be described as broadly conversational until independent testing shows:

\begin{itemize}
  \item at least 90\% mention-specific reference resolution;
  \item at least 85\% passage, causal and implication reasoning;
  \item successful twenty-turn conversations across several subjects;
  \item topic changes without loss or corruption of earlier state;
  \item reliable clarification when a question has multiple valid interpretations;
  \item unsupported-claim rate below 2\%;
  \item correction retention after one, seven and thirty days;
  \item transfer to speakers, wording and sources absent from training; and
  \item an audit proving no model or cloud route occurred in AION's response path.
\end{itemize}

These thresholds should be reported with attempt counts, source diversity and confidence
intervals.  A percentage based on a handful of attempts is not sufficient evidence.

\subsection{Minimum examination design}

The release examination should contain at least 1,000 independently graded turns distributed
across new speakers, unseen wording and source-disjoint subject matter.  No single source,
question family or template should contribute enough attempts to dominate the result.  The test
set should include:

\begin{itemize}
  \item twenty-turn conversations in science, engineering, history, literature and everyday life;
  \item conversations containing at least two deliberate topic changes;
  \item multiple same-form pronouns with different antecedents;
  \item questions with one answer, several valid answers and insufficient information;
  \item corrections that invalidate an earlier conclusion;
  \item informal and meaning-equivalent paraphrases;
  \item evidence conflicts and explicit uncertainty requirements; and
  \item delayed re-examinations after one, seven and thirty days.
\end{itemize}

\subsection{Release gate}

Broad native conversation should be declared only when all criteria pass simultaneously.  A high
grammar score cannot compensate for unsupported claims; vocabulary size cannot compensate for
state loss; and immediate success cannot compensate for failed delayed retention.  The release
receipt should report raw counts, accuracy, confidence calibration, source diversity, invalid
teacher contracts, clarification precision, unsupported claims and verified model isolation.

\section{Recommended Build Order}

The immediate implementation sequence is:

\begin{lstlisting}
1. mention-specific pronoun and reference repair
2. persistent transactional dialogue state
3. compositional arbitrary-intent parsing
4. structured knowledge retrieval from reading memory
5. native proposition and response planner
6. grammatical realiser with support maps
7. open-conversation academy and adaptive corrections
8. independent twenty-turn examinations and delayed retention
\end{lstlisting}

The reader, semantic spider, Practice Gym, Language Workshop, transfer academy and current
communication curriculum should continue in parallel.  Their role is to supply vocabulary,
evidence and supporting practice.  The majority of new engineering effort should move from
acquiring more isolated words to maintaining meaning and constructing original supported answers.

\section{Verification}

The focused and surrounding English suites cover contract filtering, object grounding, catalogue
rejection, situation models, references, correction storage, sealed teacher boundaries, stage
progression, source diversity, family caps, delayed retests, reading, practice, transfer and
online importing.

The repair was verified with 67 passing tests across:

\begin{itemize}
  \item \path{backend/tests/test_native_discourse_dialogue.py};
  \item \path{backend/tests/test_native_communication_academy.py};
  \item \path{backend/tests/test_native_english_progress.py};
  \item \path{backend/tests/test_native_reading_academy.py};
  \item \path{backend/tests/test_native_english_practice_gym.py};
  \item \path{backend/tests/test_native_semantic_transfer_academy.py}; and
  \item \path{backend/tests/test_native_online_reading_harvester.py}.
\end{itemize}

\section{Operational Checks}

After code or schema changes, operators should verify:

\begin{enumerate}
  \item the SD volume is mounted and the configured SQLite file exists;
  \item the reader, mastery curriculum and local API processes are alive;
  \item the dashboard reports a fresh snapshot with no component refresh errors;
  \item reading counters continue to advance;
  \item no invalid session contributes to capability metrics;
  \item due retests decrease or are rescheduled according to retention policy;
  \item recent dialogue contains varied families and unseen passage signatures;
  \item Qwen/student-use flags remain false for every student answer; and
  \item CPU and memory remain bounded while large retest batches run.
\end{enumerate}

Historical rows should never be deleted merely to improve a metric.  They should be assigned an
explicit validity status so the audit trail remains recoverable.

\section{Key Files and Directories}

\begin{longtable}{p{0.43\textwidth}p{0.47\textwidth}}
\toprule
Path & Purpose \\
\midrule
\path{backend/modules/aion_cognition/native_discourse_dialogue.py} & Situation models, native planning, sealed discourse exams and remediation \\
\path{backend/modules/aion_cognition/native_communication_academy.py} & Multi-turn curriculum, dialogue evidence, retention and queue controls \\
\path{backend/modules/aion_cognition/native_english_progress.py} & Consolidated evidence report \\
\path{backend/api/aion_english_intelligence_router.py} & Learning Centre APIs, progress cache and panel \\
\path{backend/modules/aion_cognition/native_reading_academy.py} & Book reading, contextual evidence and retention \\
\path{backend/modules/aion_cognition/native_semantic_transfer_academy.py} & Guided-to-free semantic transfer ladder \\
\path{backend/modules/aion_cognition/native_english_practice_gym.py} & Adaptive structured English practice \\
\path{backend/modules/aion_cognition/native_language_workshop.py} & Construction, transformation and connected-language practice \\
\path{backend/modules/aion_cognition/sqlite_connection.py} & Shared SD-safe SQLite connection policy \\
\path{backend/AION/system/aion_native_reading_service.py} & Background reading service \\
\path{backend/AION/system/aion_mastery_curriculum_service.py} & Background curriculum and communication scheduler \\
\path{backend/modules/aion_cognition/reading_library/} & Local reading curriculum and source metadata \\
\bottomrule
\end{longtable}

\section{Conclusion}

The recent work did not merely raise an English score.  It corrected what counts as evidence.
Malformed pages, incomplete answer phrases, repeated sessions and impossible teacher contracts
are now separated from genuine student performance.  AION has subsequently advanced through
short answers, explanations, paragraphs and passage discussion to its first repaired ten-turn
conversation.

The remaining problem is no longer basic vocabulary acquisition.  It is maintaining precise
connected meaning while interpreting arbitrary user intent and constructing varied, supported
responses.  The Native Conversation Runtime, mention-specific reference graph and Open
Conversation Academy are therefore the next engineering layer required to move AION from bounded
English competence toward reliable native conversation.

\end{document}
